% =====================================================================
% M3 patch: turning the refinement-thickness dichotomy into a statement
% about characteristic data, as far as the existing research allows.
%
% Sources in growupkuriyama-hub/Papers (all proofs below are taken from
% these drafts and lightly adapted to the s / \widehat s notation):
%   [A] type deletion + rare types:
%       02_fixed-h-mcfg/revision/parked_structural_results_EN.tex
%       (lem:mcfg-delete-type, prop:mcfg-rare-type-sample);
%       branch mcfg2-v79-fanout2-lower-bound-search,
%       revision/v79_fanout2_results.tex, Sec. (A)
%   [B] power transfer at fan-out two: same files, lem:mcfg-power-transfer
%   [C] no finite typing rescues the doubling family:
%       parked_structural_results_EN.tex,
%       thm:mcfg-all-typing-doubling-obstruction; v79 Sec. (C)
%   [R],[D] relaxed witnesses, abelian filters: already in main.tex
%   [F1-sl] branch mcfg2-v80-short-distribution-representatives
%       (semilattice-split abelian unions): NOT used below, unaudited.
%
% Structure of the patch
%   (0) Abstract: replace the dichotomy sentence
%   (1) Introduction: replace the paragraph on the dichotomy
%   Labels match origin/main of growupkuriyama-hub/Papers (2026-10-28).
%   (2) Section 4.2, after thm:data-from-typed-thickness:
%       Lemma type-deletion, Proposition rare-type, Corollary sandwich
%   (3) Section 4.5: replace rem:refinement-vs-data;
%       add Remark rem:fanout-one-data (cites companion Prop. 6.4)
%   (4) Section 7: new subsection "Fan-out two" (power transfer, no-lift
%       theorem), then the rewritten Open Problem and a conjecture
% =====================================================================


% ---------------------------------------------------------------------
% (0) ABSTRACT -- replace from "A further refinement-thickness
%     dichotomy ..." up to "... characteristic-data length." by:
% ---------------------------------------------------------------------
Conversely, characteristic data must contain a word of every realized
type, so the shortest word of the rarest type is an operator-independent
lower bound.  Together with a known exponential family for context-free
targets, this makes local groups exactly the finite-semigroup
pseudovarieties admitting polynomial characteristic data at fan-out one.  At fan-out two or more they remain sufficient; the
known exponential family provably fails to be substitutable there under
every finite typing, and whether local groups are necessary remains open.


% ---------------------------------------------------------------------
% (1) INTRODUCTION -- replace the paragraph beginning
%     "The obstruction of Proposition~\ref{prop:typed-thickness-exponential}
%      has an exact algebraic boundary." by:
% ---------------------------------------------------------------------
Typed thickness is not only a proof device.  Its value at the start
symbol, the length $\mu^\star(L,s)$ of the shortest word of the rarest
$s$-type of $L$, does not depend on the presentation and is a lower
bound for every set-driven learner of the class
(Proposition~\ref{prop:rare-type}), so characteristic data are
sandwiched between $\mu^\star$ and $O(\|G\|(\bar\theta_G+t_{\mathrm{typ}}))$
(Corollary~\ref{cor:typed-sandwich}).  At fan-out one, our local-group
bound combined with the exponential family of the companion CFG paper
\cite[Prop.~6.4]{kuriyamaCFG} closes the gap at the level of
pseudovarieties: polynomial characteristic data hold for every typing
with image in $\mathbf V$ if and only if $\mathbf V\subseteq\mathbf{LG}$
(Remark~\ref{rem:fanout-one-data}).  The lower-bound half is due to
\cite{kuriyamaCFG}; the contribution here is the upper bound beyond
$\mathbf{LI}\vee\mathbf{Grp}$.
At fan-out two or more, $\mathbf{LG}$ remains sufficient, and the full
refinement still has an exact boundary at $\mathbf{LG}$
(Corollary~\ref{cor:refinement-dichotomy}).  However, that family is
not $(2,s)$-substitutable under any finite typing
(Theorem~\ref{thm:doubling-no-lift}): two-component substitutability
lets powers of a factor migrate between blocks
(Lemma~\ref{lem:power-transfer}), which destroys rare types of this
kind.  Whether $\mathbf{LG}$ is necessary at fan-out two is open
(Open Problem~\ref{op:fanout-two}).


% ---------------------------------------------------------------------
% (2) SECTION 4.2 -- insert after Theorem thm:data-from-typed-thickness
%     and its trivial-typing remark
% ---------------------------------------------------------------------
\paragraph{A presentation-independent lower bound.}
For $L\subseteq\Sigma^*$ and $m\in\mathsf S$ put
$\mu_m(L):=\min\{|w|:w\in L\cap\Sigma^+,\ s(w)=m\}$ ($\min\emptyset=\infty$),
and let $\mu^\star(L,s):=\max\{\mu_m(L):m\in s(L\cap\Sigma^+)\}$ be the
length of the shortest word of the rarest realized type.  If $G$ is any
reduced good presentation of $L\cap\Sigma^+$, the typed start
nonterminals are the $S_{(m)}$ with $m\in s(L\cap\Sigma^+)$, and
$L_{S_{(m)}}=\{(w):w\in L,\ s(w)=m\}$
(Proposition~\ref{prop:output-invariants}); hence
\begin{equation}
\label{eq:root-thickness}
 \mu^\star(L,s)\le t_{\mathrm{typ}}(G,s)\qquad\text{for every such }G.
\end{equation}

\begin{lemma}[Type deletion]
\label{lem:type-deletion}
If $L\in\Cmcf{f}{s}$ and $m\in\mathsf S$, then
$L^{\setminus m}:=L\setminus s^{-1}(m)\in\Cmcf{f}{s}$.
\end{lemma}

\begin{proof}
$\widehat s^{-1}(\mathsf S^\sharp\setminus\{m\})$ is regular and
$(f,s)$-tuple-substitutable by item~(c) after
Definition~\ref{def:class-cmcf}.  Intersections preserve
$(f,s)$-tuple substitutability
(Proposition~\ref{prop:abelian-coset-filter}(i)), and $f$-MCFLs are
closed under intersection with regular languages \cite{SekiEtAl1991}.
\end{proof}

\begin{proposition}[Rare types force long data]
\label{prop:rare-type}
Let $L\in\Cmcf{f}{s}$ and $m\in s(L\cap\Sigma^+)$.
\begin{enumerate}[label=\textup{(\roman*)},leftmargin=*]
\item Every characteristic sample of $L$ for $\mathcal B^{\le r}_{f,s}$
($r\ge1$) or for $\mathcal B_{f,s}$ contains a word of type $m$.
\item The same holds for every set-driven operator that has a
characteristic sample for $L^{\setminus m}$, in particular for every
set-driven operator with characteristic samples for all of
$\Cmcf{f}{s}$.
\end{enumerate}
Hence every such characteristic sample $\mathcal F$ satisfies
$\|\mathcal F\|_+\ge\mu^\star(L,s)+1$.
\end{proposition}

\begin{proof}
(i) A finite $K\subseteq L$ without a word of type $m$ lies in
$L^{\setminus m}\in\Cmcf{f}{s}$ (Lemma~\ref{lem:type-deletion}), so
Proposition~\ref{prop:reconstruction-soundness} applied to the target
$L^{\setminus m}$ gives
$L(\mathcal B(K))\subseteq L^{\setminus m}\subsetneq L$.
(ii) If $\mathcal F$ is characteristic for $L$ and avoids type $m$, and
$\mathcal F'$ is characteristic for $L^{\setminus m}$, then
$\mathcal F\cup\mathcal F'$ lies between $\mathcal F$ and $L$ and between
$\mathcal F'$ and $L^{\setminus m}$, so the operator would have to return
a grammar for both $L$ and $L^{\setminus m}\ne L$.
\end{proof}

\begin{corollary}[Sandwich]
\label{cor:typed-sandwich}
Let $L\in\Cmcf{f}{s}$ and let $G$ be a reduced good fan-out-$\le f$,
rank-$\le r_{\max}$ presentation of $L\cap\Sigma^+$.  The minimum encoded
length $\mathrm{cs}(L)$ of a characteristic sample of $L$ for
$\mathcal B^{\le r_{\max}}_{f,s}$ satisfies
\[
 \mu^\star(L,s)+1\;\le\;\mathrm{cs}(L)\;\le\;
 N_{\mathrm{typ}}\bigl(1+2\bar\theta_G+r_{\max}t_{\mathrm{typ}}(G,s)\bigr)+1 .
\]
The lower bound is independent of $G$ and holds for every set-driven
operator with characteristic samples for all of $\Cmcf{f}{s}$.
\end{corollary}

\begin{proof}
Proposition~\ref{prop:rare-type} and
Theorem~\ref{thm:data-from-typed-thickness}; the final $+1$ accounts for
$\eps$.
\end{proof}

% One sentence after the corollary:
% "Thus the full refinement can be exponentially thicker than necessary
%  only below the start symbol: by \eqref{eq:root-thickness} its
%  start-level part is a genuine lower bound."


% ---------------------------------------------------------------------
% (3a) SECTION 4.5 -- replace Remark rem:refinement-vs-data by:
% ---------------------------------------------------------------------
\begin{remark}[Where the exponential sits]
\label{rem:refinement-vs-data}
The witnesses of Corollary~\ref{cor:refinement-dichotomy} present
$\Sigma^+$.  Every element of $\mathsf S$ is the type of a word of length
at most $|\mathsf S|$ (shorten any longer word at a repeated prefix
product), so $\mu^\star(\Sigma^+,s)\le|\mathsf S|$, whereas
$t_{\mathrm{typ}}(G_n,s)\ge2^n$.  The exponential thus lies strictly
below the start symbol, where
Corollary~\ref{cor:typed-sandwich} provides no lower bound.  A
data lower bound needs instead a target whose \emph{start-level}
thickness $\mu^\star$ is large while $\|G\|\bar\theta_G$ is small.  At
fan-out one such targets exist (Remark~\ref{rem:fanout-one-data});
at fan-out two this is Open Problem~\ref{op:fanout-two}.
\end{remark}


% ---------------------------------------------------------------------
% (3b) SECTION 4.5 -- insert after cor:refinement-dichotomy
%      (cites [kuriyamaCFG, Prop. 6.4] as an external result)
% ---------------------------------------------------------------------
\begin{remark}[Fan-out one: the boundary for data length]
\label{rem:fanout-one-data}
At fan-out one the boundary of
Corollary~\ref{cor:refinement-dichotomy} is also the boundary for
characteristic-data length, at the level of pseudovarieties.  If
$\mathbf V\subseteq\mathbf{LG}$, Theorem~\ref{thm:local-group-polydata}
with $f=1$ gives polynomial characteristic data for every typing with
image in $\mathbf V$.  If $\mathbf V\not\subseteq\mathbf{LG}$, then
$U_1=\{1,\zeta\}\in\mathbf V$ (Theorem~\ref{thm:lg-characterization}), and
the companion paper supplies the exponential family
\cite[Prop.~6.4]{kuriyamaCFG}: over
$\Gamma_{\mathrm{br}}=\{\mathtt a,\lfloor,\rfloor,\lceil,\rceil\}$, with
$s_U$ sending $\lceil,\rceil$ to $\zeta$ and the other letters to $1$,
the CFGs $S_0\to E_n$, $E_0\to\mathtt a$,
$E_i\to\lfloor E_{i-1}E_{i-1}\rfloor\mid\lceil E_{i-1}\rceil$ are reduced,
good, of rank two, with $\|G_n\|=O(n)$ and $\bar\theta_{G_n}=O(n)$, and
generate Clark--Eyraud substitutable languages $L_n$, hence
$L_n\in\Cmcf{1}{s_U}$.  The only word of $L_n$ of type $1$ has length
$3\cdot2^n-2$, so Proposition~\ref{prop:rare-type} forces every
characteristic sample, for $\mathcal B^{\le2}_{1,s_U}$ and for every
set-driven learner of $\Cmcf{1}{s_U}$, to have encoded length at least
$3\cdot2^n-1$.  We use this family only here and only through the cited
result; whether every single typing outside $\mathbf{LG}$ admits such a
family over its own alphabet is open.
\end{remark}


% ---------------------------------------------------------------------
% (4) SECTION 7 (Conclusion) -- insert before the Open Problem a short
%     subsection, then replace the Open Problem
% ---------------------------------------------------------------------
\subsection*{Fan-out two: why the fan-out-one lower bound does not lift}

\begin{lemma}[Power transfer]
\label{lem:power-transfer}
Let $f\ge2$, $L\in\Cmcf{f}{s}$, $u\in\Sigma^+$, and let $i_0,\pi\ge1$ be
the index and period of $s(u)$, so $s(u)^{i+\pi}=s(u)^i$ for $i\ge i_0$.
Let $a,c\ge i_0$ and $d\in\pi\mathbb Z_{>0}$.  If some word of $L$ has the
factor $u^{a+d+c+2}$, then for all $X,Z\in\Sigma^*$ and $Y\in\Sigma^+$
\[
 Xu^{a}Yu^{d+c}Z\in L\iff Xu^{a+d}Yu^{c}Z\in L .
\]
\end{lemma}

\begin{proof}
Write the given word as $Pu^{a+d+c+2}Q$.  The pairs
$\tuple x=(u^a,u^{d+c})$ and $\tuple y=(u^{a+d},u^c)$ have equal
componentwise $s$-types, and the admissible context
$E=Pu\,\blank_1u\blank_2Q$ gives
$E[\tuple x]=E[\tuple y]=Pu^{a+d+c+2}Q\in L$.  Hence
$\D^{(2)}_L(\tuple x)=\D^{(2)}_L(\tuple y)$; test both in
$X\blank_1Y\blank_2Z$, which is admissible because $Y\ne\eps$.
\end{proof}

\begin{theorem}[The doubling family does not lift]
\label{thm:doubling-no-lift}
Let $s:\Gamma_{\mathrm{br}}^+\twoheadrightarrow\mathsf S$ be any typing
and $L_n$ as in
Remark~\ref{rem:fanout-one-data}.  If $n\ge2|\mathsf S|+3$,
then $L_n\notin\Cmcf{f}{s}$ for every $f\ge2$.
\end{theorem}

\begin{proof}
Let $i_0,\pi$ be the index and period of $s(\lfloor)$; then
$i_0+\pi\le|\mathsf S|+1$, so $n\ge2i_0+\pi+2$.  Put $a=i_0$, $d=\pi$,
$c=n-a-d-2\ge i_0$.  The ceiling-free word
$W_n=\lfloor W_{n-1}W_{n-1}\rfloor$ starts with $\lfloor^n$.  Write
$W_n=X\lfloor^aY\lfloor^{d+c}Z$ with $X=\lfloor$, with $\lfloor^a$ the next
$a$ letters of the leading run, and with $\lfloor^{d+c}$ the last $d+c$
letters of the leading run $\lfloor^{n-1}$ of the second copy of
$W_{n-1}$; then $Y$ contains the first $\mathtt a$.  If
$L_n\in\Cmcf{2}{s}$, Lemma~\ref{lem:power-transfer} with $u=\lfloor$
puts $X\lfloor^{a+d}Y\lfloor^cZ$ into $L_n$.  This word is ceiling-free
but begins with $\lfloor^{n+d}\mathtt a$, whereas the only ceiling-free
word of $L_n$ begins with $\lfloor^n\mathtt a$.  The case $f>2$ follows
because every $(f,s)$ condition contains arity two.
\end{proof}

\begin{remark}
The transfers of Lemma~\ref{lem:power-transfer} move material without
changing length (the two words have equal Parikh vectors), so they do
not by themselves shorten representatives.  They show instead that a
long run inside one word of a fan-out-two target forces many other words
into the target, which is what a rare type must avoid.
\end{remark}

\begin{openproblem}[Characteristic-data length at higher fan-out]
\label{op:fanout-two}
Fix $f\ge2$ and $r_{\max}\ge2$.  For a typing $s$ with
$\mathsf S\notin\mathbf{LG}$, are there targets $L_n\in\Cmcf{f}{s}$ with
reduced good presentations $G_n$ such that
$\|G_n\|\bar\theta_{G_n}=n^{O(1)}$ but every characteristic sample for
$\mathcal B^{\le r_{\max}}_{f,s}$ has length $n^{\omega(1)}$?  By
Proposition~\ref{prop:rare-type} it suffices that
$\mu^\star(L_n,s)=n^{\omega(1)}$, and the lower bound then holds for every
set-driven learner of $\Cmcf{f}{s}$.
\end{openproblem}

\begin{conjecture}[Rare types are short at fan-out two]
\label{conj:rare-types}
For all $\Sigma$, $s$ and $r_{\max}$ there is a polynomial $p$ with
$\mu^\star(L,s)\le p(\|G\|,\bar\theta_G)$ whenever $L\in\Cmcf{2}{s}$ and
$G$ is a reduced good rank-$\le r_{\max}$ presentation of $L\cap\Sigma^+$.
\end{conjecture}

% One paragraph after the conjecture:
% "The conjecture holds when $\mathsf S\in\mathbf{LG}$, by
%  \eqref{eq:root-thickness} and the bound on $t_{\mathrm{typ}}$ in the
%  proof of Theorem~\ref{thm:local-group-polydata},
%  and for abelian-filtered targets under every typing, because the
%  characteristic sample of Theorem~\ref{thm:mcfg-abelian-filter-data}
%  contains a word of each realized type (Proposition~\ref{prop:rare-type}).
%  It fails at fan-out one (Remark~\ref{rem:fanout-one-data}).
%  If it holds, the rare-type method cannot settle
%  Open Problem~\ref{op:fanout-two}."


% =====================================================================
% NOTES FOR THE AUTHOR (not for the manuscript)
%
% * rem:fanout-one-data cites [kuriyamaCFG, Prop. 6.4] as an external
%   result and credits the lower-bound half to it; no proof is
%   reproduced, so there is no duplicate claim.  Replace the sentence
%   "No lower bound or separating family from the related CFG manuscript
%   is used" (Introduction, Conclusion) by "The fan-out-one lower bound of
%   Remark~\ref{rem:fanout-one-data} is cited from \cite{kuriyamaCFG};
%   no other result of that manuscript is used."  The bibliography note
%   on Prop. 6.4 then becomes consistent (M7).  Make the version
%   containing Prop. 6.4 publicly available (e.g. arXiv) before I&C
%   submission so that referees can check it.
% * thm:doubling-no-lift needs only s(\lfloor); its threshold 2|S|+3 is
%   cruder than the window-specific one (n >= max{3,k+l+1}) in
%   parked_structural_results_EN.tex, prop:fixed-window-doubling-fanout-two.
% * Status: [A],[B],[C] were drafted 2026-10-18 for independent audit and
%   have not been merged into main; check them again before adoption.
% =====================================================================
